\section{The XY flow, admissible tuning, and spin-field limit}
\label{xyf:section}

We apply the representation calculus of Proposition~\ref{xyh:calculus}
to the ordinary nearest-neighbor XY model.  Throughout this section,
\(b\) denotes spin inverse temperature,
\[
 \mu=1+s_{11},\qquad d_c=1+s_{12},\qquad s=s_{22},
 \qquad s\ge d_*s_{12}^2.
\]
The constant \(d_*>0\) will be chosen below.  The interaction includes
the compensating factor \(e^{+U(w,p)}\), with
\[
 U(w,p)=\frac12\sum_e
  \bigl(s_{11}w_e^2+2i s_{12}w_ep_e+s_{22}p_e^2\bigr),
 \qquad p=\partial\psi.
\]
The covariance prescription and the regularized interpretation at
\(s=s_{12}=0\) are those of Proposition~\ref{xyh:compensation}.
In particular, the prescription is not integration against a positive
joint law of \((w,\psi)\).  Positivity will be used only for the completely
restored vortex/transverse ensemble of Proposition~\ref{xyh:branch}.

\subsection{The exact recurrence}

For a full block \(B\), write
\begin{equation}
 \mathcal E_B(t)=\frac12\sum_e\mathrm{wt}_B(e)
 \bigl(t_{11}w_e^2+2i t_{12}w_e(\partial\psi)_e
                  +t_{22}(\partial\psi)_e^2\bigr).
 \label{xyf:tensor}
\end{equation}
Here \(\mathrm{wt}_B\) assigns half of an edge to each endpoint in \(B\).
Thus the weights sum to one when the blocks partition the whole grid.
The tensor norm is the maximum of the three coefficient absolute values.
The size of a remainder is the supremum over its supports of
\(A^{|X|}\) times its positive representation size.  Representation size
means the recursive sum in Proposition~\ref{xyh:calculus}, with each term
weighted by its live-root weight and the product of its opaque completed
coefficient sizes;
it is not a pointwise supremum over future complex Gaussian fields.

We record precisely which parts of that proposition enter the flow.
They hold uniformly in the number of integrated scales and in the
finite-grid size.  A microscopic \(F_b\) factor has size \(C/b\), a
nonzero charge has a summable bound \(e^{-cbq^2}\), and the nonlinear
part of a microscopic compensation factor has size \(C|s_\bullet|^2\).
A tensor singleton has size \(C|t|\).  Concatenating compatible histories
multiplies their sizes, up to a fixed constant per input.  On a small
\emph{linear} input, neutral localization gives the following gains:
\[
 \begin{array}{c|c|c}
 \text{class}&\text{subtracted polynomial}&\text{remainder gain}\\ \hline
 \text{bulk, even}&\text{constant and quadratic}&CL^{-3}\\
 \text{wall, even}&\text{constant}&CL^{-2}\\
 \text{defect}&\text{neutral constant}&CL^{-1}.
 \end{array}
\]
Every small charged linear input instead has gain \(CL^{-4}\).
The extracted coefficients, in normalized fine-scale directions, are
bounded by \(C\) times the input size.  The same history estimates bound
terminal evaluations and fixed smooth-source derivatives.  Finally,
local identities are exact, their finite-depth values are continuous in
the admissible parameters including the degenerate boundary, and their
field arguments are analytic in the source neighborhoods used below.

\begin{proposition}[Bulk recurrence]\label{xyf:recurrence}
For every sufficiently small fixed \(\theta>0\), the geometric
parameters, then \(L\), and then \(A\) can be chosen so that the exact
bulk activities have the decomposition
\[
 K_j(X)=\ind_{X=\{B\}}\mathcal E_B(t_j)+\mathcal P_j(X),
 \qquad t_0=(s_{11},s_{12},s_{22}).
\]
If \(r_j\) is their remainder size, then, while the quantities on the
right are small,
\begin{align}
 |t_{j+1}-t_j|_\infty&\le C r_j,\label{xyf:increment}\\
 r_{j+1}&\le\theta r_j+C_{L,A}(|t_j|_\infty+r_j)^2,
       \label{xyf:remainder}\\
 r_0&\le C_{L,A}\bigl(b^{-1}+|t_0|_\infty^2\bigr).
       \label{xyf:initial}
\end{align}
At every fixed depth these coefficients and the needed fixed field
derivatives are continuous on the set where the preceding bounds are
retained.  The recurrence also holds, with the same bulk tensors,
on an auxiliary torus of side \(F L^j\), through scale \(L^j\), if
the fixed integer \(F\) is sufficiently large.
\end{proposition}

\begin{proof}
Use the temporary-component, relocation, and constant-decoration
identities of Proposition~\ref{rg:partition-map}.  They are algebraic
identities for local activities, used with the enlarged fixed separation
radii proved in Lemma~\ref{xyh:history}. Their hypotheses here are finite
shell range, the retained temporary exclusions, one-step locality of
assigned pieces, and multiplicative bounds for compatible inputs.
Finite range follows from Lemma~\ref{xyh:kernels}; the history
construction supplies the other hypotheses, including the fact that
evaluated coefficients are scalar factors and do not retain their
extinct microscopic leaves.

For completeness, the relevant counts are as follows.  The maximum
fine-block count \(n_0\) of a small component and the connectivity radius
are fixed before \(L\).  The number of small animals containing a
specified fine block is therefore a constant.  A coarse block has
\(O(L^2)\) fine blocks, giving at most \(CL^2\) small linear inputs per
coarse output.  Their neutral and charged contributions have total
operator bound
\begin{equation}
 C L^2(L^{-3}+L^{-4})\le C/L.
 \label{xyf:bulkcount}
\end{equation}
In the neutral bulk quadratic remainder, replace each argument by its
affine approximation.  The two-factor difference is the sum of one
term with the first argument replaced and one with the second replaced.
The error direction has size \(C/L^2\) and the other direction size
\(C/L\), so each term has size \(C/L^3\).  The first-order term
vanishes because the full neutral projection is even under
\((w,\psi)\mapsto(-w,-\psi)\).  This explains the gain used in
\eqref{xyf:bulkcount} without assigning a Taylor gain to a nonlinear
convolution.

For a large input, the closure estimate leaves an exponential weight
gain at least \(A^{-|X|/2}\), after the fixed per-input history cost.
Taking \(A\) large after \(L\) makes the sum of its linear contributions
as small as prescribed.  To check nonlinear terms, root a compatible
family at a specified output block.  If it has \(q\) fine components
and a total of \(v\) fine blocks, a depth-first traversal of the fine
animals and a tree connecting their closures gives at most
\[
 C^v(CL^2)^{2q}
\]
choices.  The long jumps of that tree have range \(CL\) in fine-block
units.  For large components the weight gain sums these choices; the
total size of the small components is at most \(qn_0\).  Thus all
terms with \(q\ge2\) sum to
\(C_{L,A}(|t_j|+r_j)^2\).  Use a temporary output weight \(4A\) before
the constant-decoration step.  If \(a_B\) is the assigned constant,
then \(|a_B|\le C_{L,A}(|t_j|+r_j)\) and
\[
 d_B=(1+a_B)^{-1}-1=-a_B+O(a_B^2).
\]
Overlapping a decoration with an old occupied site costs at most two
choices, paid by the extra weight.  A decoration at a hole supplies a
factor \(|d_B|\).  The subset sum converges when
\(A\sup_B|d_B|<1/8\).  Every term containing an old component and a
decoration is quadratic; pure decorations are quadratic apart from
the displayed cancelling linear singleton.  Relocated multiple-object
terms are also quadratic, because their original temporary exclusions
are retained.  These counts prove the stated nonlinear bound.

The entire expectation of a tensor singleton is assigned to the coarse
singleton.  Its nonconstant part is exactly the same quadratic form;
its covariance contraction is a bounded constant.  At any specified
\emph{fine} block only a bounded number of small supports contribute
quadratic coefficients, each bounded by \(Cr_j\).  Averaging over
orientation signs and anchors and summing all fine blocks in a full
coarse block preserves translation and square symmetries.  Hence the
sum is precisely \(\mathcal E_B(t_{j+1}-t_j)\), with coefficient
bound \(Cr_j\).  There is no factor \(L^2\) in this coefficient:
the sum over fine sites is the edge sum already present in
\eqref{xyf:tensor}.  Complex conjugation combined with
\(\psi\mapsto-\psi\) makes the three extracted parameters real.
This proves \eqref{xyf:increment} and \eqref{xyf:remainder}.

At the first independent piece, the linear compensation factors give
the stated tensor and their designated covariance constants.
The remaining single primitives cost \(C/b\),
\(\sum_{q\ne0}e^{-cbq^2}\), or \(C|t_0|^2\).  Each edge has two
endpoints and only finitely many choices of primitive type.  The same
animal and decoration counts therefore give \eqref{xyf:initial}.

All these identities may first be performed with a finite Fourier
cutoff and strict covariance regularization.  At a fixed depth there
are finitely many geometric choices for each coefficient.  The charge
sums retain Gaussian domination; the contour and Taylor parameters
range over compact sets.  The completed finite-dimensional integrals
are continuous, including singular real Gaussian covariances, and
the scalar denominators stay away from zero.  The domination in
Proposition~\ref{xyh:calculus} consequently permits removing the
regularizations and taking parameter limits in each of these exact
identities.  It also permits fixed smooth-source Cauchy derivatives.
This proves the asserted continuity without a positivity assumption on
the mixed Gaussian pair.

For the torus take \(F\) larger than the fixed support, kernel, and
assignment clearances.  Every small input through the indicated scale
has an unambiguous lift to the plane, with its required neighborhood;
a winding input is large.  The local tensor assignments are therefore
identical to the plane assignments.  Large winding inputs obey the
large-support weight estimate.  The transverse harmonic directions
are added at the terminal integration, as prescribed in
Lemma~\ref{xyh:kernels}, and do not change any earlier local map.
\end{proof}

\subsection{The lower-face diagnostic}

The restriction \(s\ge d_*s_{12}^2\) is essential to the history bounds,
so a three-dimensional shooting argument must respect it.  We derive
the required boundary information from a positive finite ensemble.
On the auxiliary torus let
\[
 \mathrm d\nu(q,X)=Z^{-1}e^{-\|V_q\|^2/2}
     \prod_e f_b((X+V_q)_e)\,\mathrm d\gamma_{P_T}(X),
 \quad V_q=2\pi\sqrt b\,\partial Gq,
 \quad \sum_vq_v=0.
\]
Here \(P_T\) includes the harmonic transverse directions.  This is
a positive probability measure.  It is used only as the diagnostic
ensemble; no identification with a periodic XY spin law is required.

Choose a real torus function \(\varphi\) for which
\(g=\partial\varphi\) has norm one and is a lowest nonconstant
frequency mode.  For side \(M\), one may normalize a sine wave:
\(\|\varphi\|_\infty\le C\) and
\(\|\nabla^r\varphi\|_\infty\le C_rM^{-r}\).
Put
\[
 a=s/d_c,\qquad B_c=d_c+\mu a,
 \qquad w_{\rm ext}=(\lambda+i a\tau)g,
 \qquad \psi_{\rm ext}=\tau\varphi.
\]
Let \(Z(\lambda,\tau)\) be the fully integrated compensated partition
function evaluated at these external fields, with the same scalar
normalization as \(Z(0,0)\).  Denote the Hessian at zero of
\(\log[Z(\lambda,\tau)/Z(0,0)]\) by \(H\).

\begin{lemma}[Exact diagnostic Hessian]\label{xyf:hessian}
Write \(w=X+V_q\), \(V_g=(V_q,g)\), \(f(w)=\prod_e f_b(w_e)\), and
\[
 O=(\mu-d_c)V_g+d_c\partial_g\log f(w).
\]
Then
\begin{align}
 H_{\tau\tau}&=s+\mu a^2-B_c^2\Var_\nu(V_g),
       \label{xyf:htt}\\
 H_{\lambda\tau}/i
     &=B_c\bigl(s_{12}+\Cov_\nu(O,V_g)\bigr),
       \label{xyf:hlt}\\
 H_{\lambda\lambda}
     &=-\mu-d_c^2+2\mu d_c
        +d_c^2\E_\nu\partial_g^2\log f(w)+\Var_\nu(O).
       \label{xyf:hll}
\end{align}
If \(|H_{\lambda\lambda}|\le1\) and the parameters lie in a fixed
sufficiently small neighborhood of zero, then \(\Var_\nu(O)\le C\),
with \(C\) independent of the torus size and \(b\) in the asserted
low-temperature range.  Consequently, if additionally
\(|H_{\lambda\tau}|\le C_1\delta\), then
\begin{equation}
 \Var_\nu(V_g)\ge c_0(|s_{12}|-C_2\delta)_+^2.
 \label{xyf:variance}
\end{equation}
\end{lemma}

\begin{proof}
Translate the full fields before undoing the compensation.  The
longitudinal vortex argument becomes \(V_q+d_c\lambda g\): the
\(\tau\)-terms cancel because \(s-d_ca=0\).  The linear source on
\(w\) is \((\mu\lambda+iB_c\tau)g\), the source on \(p\) is
\(id_c\lambda g\), and the direct translation factor is
\[
 \exp\{-\mu\lambda^2/2-iB_c\lambda\tau
                  +(s+\mu a^2)\tau^2/2\}.
\]
The constant scalar mode enforces \(\sum q_v=0\).
At each such charge configuration the \(p\)-source gives the further
factor \(\exp\{-d_c\lambda V_g-d_c^2\lambda^2/2\}\).
Since \((X,g)=0\), the linear \(w\)-source is evaluated at
\(V_g+d_c\lambda\).  Combining these factors gives the exact ratio
\begin{align}
 \frac{Z(\lambda,\tau)}{Z(0,0)}
 &=\exp\!\left\{\frac{-\mu-d_c^2+2\mu d_c}{2}\lambda^2
        +iB_c(d_c-1)\lambda\tau
        +\frac{s+\mu a^2}{2}\tau^2\right\}\notag\\[-2pt]
 &\quad\times\E_\nu\left[
    e^{((\mu-d_c)\lambda+iB_c\tau)V_g}
    \frac{f(w+d_c\lambda g)}{f(w)}\right].
 \label{xyf:exactratio}
\end{align}
The identity first holds in the strict regularizations and then passes
to the prescribed finite-graph limit by Proposition~\ref{xyh:compensation}.
Real positivity of \(f\) makes the displayed logarithmic derivatives
well defined.  Differentiating \eqref{xyf:exactratio} proves all three
Hessian identities; in particular the direct mixed term is
\(iB_c(d_c-1)=iB_cs_{12}\).

The branch estimate in Lemma~\ref{xyh:branch-bounds} gives
\((\log f_b)''\ge-C_f\).  As \(\sum_e g_e^2=1\),
\[
 \partial_g^2\log f(w)
     =\sum_e g_e^2(\log f_b)''(w_e)\ge-C_f.
\]
Rearranging \eqref{xyf:hll}, and bounding its deterministic terms,
gives \(\Var_\nu(O)\le C\).  Because \(B_c\) stays bounded away
from zero, \eqref{xyf:hlt} yields
\(
 |\Cov_\nu(O,V_g)|\ge(|s_{12}|-C_2\delta)_+.
\)
Cauchy--Schwarz for the positive measure \(\nu\) now gives
\eqref{xyf:variance}.
\end{proof}

\begin{proposition}[Admissible infinite trajectory]\label{xyf:shooting}
For all sufficiently large \(b\), there is a choice
\[
 |s_{11}|+|s_{12}|+s_{22}=O(b^{-1}),\qquad
 s_{22}\ge d_*s_{12}^2,
\]
for which the exact bulk flow satisfies
\begin{equation}
 |t_j|_\infty\le T\eps\rho^j,
 \qquad r_j\le2\eps\rho^j,
 \qquad \eps=C_*/b,
 \quad 0<\theta<\rho<1.
 \label{xyf:survivor}
\end{equation}
The constants and the lower bound on \(b\) do not depend on a
correlation order or a finite-volume cutoff.
\end{proposition}

\begin{proof}
First choose \(d_*\) small using the universal variance constant in
Lemma~\ref{xyf:hessian}.  That constant requires only
\(|H_{\lambda\lambda}|\le1\) and a fixed small parameter
neighborhood; it does not require a decay estimate or a choice of
\(T\).  Fix the history geometry and a sufficiently small
\(\theta>0\).  Choose \(L\) large enough for the bulk, wall, and
defect small-input contraction counts in
Propositions~\ref{xyf:recurrence} and~\ref{xyf:forced}.
Next choose the defect weight \(A_D\) sufficiently large for the
large-input contraction and decoration estimates, and then choose
the untwisted weight \(A_U\) sufficiently large, with
\(A_U/A_D\) large enough for the support and assignment counts accompanying
the same-representation transport comparison in
Proposition~\ref{xyf:forced}.  In the bulk recurrence and in this
shooting argument set \(A=A_U\).  Both weights are now fixed and
will not be enlarged after a survivor is chosen.  The estimates
requiring these choices depend only on the history geometry and
the uniform one- or two-source bounds, also valid in the separated
many-source decaying regime; they do not depend on a correlation
order or on the eventual surviving parameters.
Fix \(\rho\in(\theta,1)\), and fix
\(\rho<\rho_1<\rho_2<1\) for the forced recurrences, together with
the diagnostic torus factor \(F\).  All recurrence, translation,
and terminal constants below are evaluated at these fixed choices.
Choose \(C_*\) to dominate the microscopic \(b^{-1}\) constant.
Choose \(T\) large afterwards, and finally choose the lower bound
on \(b\) large enough for the bulk shooting estimates and the
wall and defect smallness conditions with these same weights.

Let \(h=T\eps\), \(\delta_j=h\rho^j\), and use the parameter domain
\begin{equation}
 P_h=\{(s_{11},s_{12},s): |s_{11}|,|s_{12}|\le h,
                     \ d_*s_{12}^2\le s\le h\}.
 \label{xyf:parameterdomain}
\end{equation}
It is a closed topological ball when \(d_*h<1\).  Its lower face
is the graph \(s=d_*s_{12}^2\).  If
\(|t_i|_\infty\le\delta_i\) through step \(j-1\), induction in
\eqref{xyf:remainder} gives \(r_j\le2\eps\rho^j\): it suffices that
\[
 C_{L,A}(T+2)^2\eps\le2(\rho-\theta),
\]
and \eqref{xyf:initial} gives the initial inequality by increasing
\(b\).  At a coordinate equality
\(|t_{j,\alpha}|=\delta_j\), the increment estimate gives
\[
 |t_{j+1,\alpha}|\ge\delta_j-2C\eps\rho^j>\delta_{j+1}
\]
on the same signed side, provided \(T(1-\rho)>2C\).
This strict crossing will control boundary contacts of the exit map.

We next exclude an exit through the positive \(22\)-face for a
parameter on the lower face of \(P_h\).  Suppose the trajectory
has stayed inside through step \(j-1\), and at step \(j\) has
\(t_{22,j}\ge\delta_j\).  The increment bound implies
\(|t_j|_\infty\le M_0\delta_j\), with fixed \(M_0\), and the
remainder bound still holds at step \(j\).  There is no such event
at \(j=0\), since \(d_*h^2<h\).

Perform the calculation on the auxiliary torus of side \(FL^j\).
After removing its designated constants there are at most \(C_F\)
terminal blocks.  The terminal expression is its tensor contribution
plus an error whose smooth-source derivatives of order at most two
are bounded by
\(C_F(r_j+|t_j|^2)\).  Indeed a single tensor expectation has its
unmodified external-field quadratic and a field-independent covariance
constant; every other linear input is a remainder, and every
multi-input term has at least two tensor/remainder factors.  The
terminal history estimate bounds their integrals and Cauchy
derivatives on fixed small circles in \((\lambda,\tau)\).
The normalized zero-source partition is \(1+O_F(|t_j|+r_j)\),
so logarithmic differentiation has the same error bound.  The chosen
sine mode obeys all the required box-scale test bounds uniformly in
\(j\). On the fixed complex \((\lambda,\tau)\)-polydisc, its scalar
background, including the constant value on each terminal block,
contributes at most \(\exp\{C h_{L^j}(k;X)\}\) for the current charge
array \(k\) on a terminal support \(X\). This cost is paid by the
final normalized-test allowance in Lemma~\ref{xyh:budget}; the test
class defining \(h_m\) includes its zeroth-order bound and does not
quotient out constants. The finitely many terminal test events only
change \(C_F\), which is fixed before \(T\). The uniform majorants
on smaller source polydiscs in
Proposition~\ref{xyh:calculus}(iv)--(v) justify the Cauchy derivatives
also at the degenerate endpoint \(s=0\).
Consequently
\begin{align}
 H_{\lambda\lambda}&=t_{11,j}
                  +O_F(r_j+|t_j|^2),\notag\\
 H_{\lambda\tau}/i&=t_{12,j}+a t_{11,j}
                  +O_F(r_j+|t_j|^2),\notag\\
 H_{\tau\tau}&=t_{22,j}-2a t_{12,j}-a^2t_{11,j}
                  +O_F(r_j+|t_j|^2).
 \label{xyf:terminalhessian}
\end{align}
All constants here are independent of \(T\).

For small \(h\), the first line ensures
\(|H_{\lambda\lambda}|\le1\); the second gives
\(|H_{\lambda\tau}|\le C\delta_j\).  Equations
\eqref{xyf:htt} and \eqref{xyf:variance} therefore imply, on the
lower face,
\begin{equation}
 H_{\tau\tau}\le d_*s_{12}^2
       +C d_*^2s_{12}^4
       -c_1(|s_{12}|-C_2\delta_j)_+^2
       \le C_3\delta_j^2.
 \label{xyf:bottomupper}
\end{equation}
For the last inequality use
\((x-y)_+^2\ge x^2/2-y^2\), choose \(d_*<c_1/4\), and then
make \(h\) small enough to absorb the fourth-order term.
On the other hand, \(|a|\le Ch\) and the last line of
\eqref{xyf:terminalhessian} imply
\[
 H_{\tau\tau}\ge\delta_j-C h\delta_j
                 -C_F\bigl(2\delta_j/T+M_0^2\delta_j^2\bigr)
             \ge\tfrac12\delta_j.
\]
Choose \(T\) sufficiently large, then \(h\) sufficiently small.
This contradicts \eqref{xyf:bottomupper}.  These choices are not
circular: the variance bound preceded \(d_*\), and the terminal
constant \(C_F\) was fixed before \(T\).

Here is the topological conclusion, including boundary intersections.
Put \(x_j=t_j/\delta_j\), and interpolate consecutive \(x_j\)'s
linearly in a time variable.  If no infinite survivor exists, every
parameter has a first contact with the boundary of the cube
\(Q=[-1,1]^3\).  Denote that point by \(E(s_\bullet)\).
This defines a continuous map \(P_h\to\partial Q\), as follows.
Put \(D_j=\{s_\bullet:x_i(s_\bullet)\in Q\text{ for }i<j\}\).
Finite-recursion continuity on these relative domains makes the
\(D_j\) nested compact sets.  If their intersection is empty, then
\(D_N=\varnothing\) for some finite \(N\).
Let \(n\) be the first index at which \(x_n\notin Q\).
Strict crossing implies that every \(x_i\), \(i<n-1\), is in
the interior; the first-contact point lies on the final segment,
possibly at its initial endpoint.  Along any convergent parameter
sequence, pass to a subsequence with the same \(n\) and with
convergent first-contact interpolation fractions.  Closedness of
\(D_n\) and relative continuity give convergence of both segment
endpoints.  Their limiting interpolated point lies on \(\partial Q\).
There can be no earlier boundary contact: an equality at an earlier
integer time would force the next endpoint outside \(Q\), contrary
to membership in \(D_n\).  If the final segment starts on the boundary,
strict outward crossing makes that initial point its unique boundary
point; otherwise it starts in the interior and again has a unique
boundary point.  This also covers a final endpoint exactly on the
boundary.  Thus every subsequential limit is the first-contact point
at the limiting parameter, proving continuity including at corners.

On the top and on either vertical face of \(P_h\), the initial
point already belongs to \(\partial Q\), so
\(E(s_\bullet)=s_\bullet/h\).  The lower-face exclusion just proved
implies that its first-contact image avoids the relative interior
of the top cube face.  Indeed a contact there would require the next
discrete endpoint to have \(t_{22,j}\ge\delta_j\), while every
earlier endpoint was retained.  At the intersections with a vertical
face the initial image lies on that vertical face; it is never in
the relative interior of the top face.

The boundary map \(E|_{\partial P_h}\) has degree one.  For example,
the point \((0,0,1)\) has exactly the preimage \((0,0,h)\), in the
relative interior of the top parameter face; in a neighborhood there
the boundary map is the orientation-preserving map
\((s_{11},s_{12})\mapsto(s_{11}/h,s_{12}/h)\).
There are no preimages on the lower or vertical faces.  But any map
from the ball \(P_h\) to the sphere \(\partial Q\) restricts to a
null-homotopic boundary map, by contracting the ball to a point, and
has boundary degree zero.  This contradiction proves the existence
of a survivor.  Fix one such parameter choice for each \(b\).
\end{proof}

\subsection{Walls, sources, and the single amplitude}

The block construction is available for every integer side length,
including partial terminal blocks, as in Lemma~\ref{rg:geometry}.
Fix the surviving bulk parameters, independently of the finite grid
and of translations of the nested cuts.  Write
\(u_j=|t_j|_\infty+r_j\le C\eps\rho^j\).
Use the weights \(A_U=A\) and \(A_D\) fixed before shooting in
Proposition~\ref{xyf:shooting}.  The bulk and untwisted wall sizes
are measured at \(A_U\), and the source-defect sizes at \(A_D\).
The same-representation comparison of Proposition~\ref{xyh:calculus}
handles bounded early translations before relocation. The fixed ratio
\(A_U/A_D\) pays only the accompanying support and assignment counts.

\begin{proposition}[Forced recurrences]\label{xyf:forced}
For untwisted wall activities let \(w_j\) denote their full size.
For twisted defect activities let \(v_j\) denote their full size.
There are fixed \(\rho<\rho_1<\rho_2<1\) such that
\begin{align}
 w_{j+1}&\le\theta w_j+C u_j+C(u_j+w_j)^2,
             &w_j&\le C\eps\rho_1^j,
             \label{xyf:wallrec}\\
 v_{j+1}&\le\theta v_j+C(u_j+w_j)+C(u_j+w_j+v_j)^2,
             &v_j&\le C\eps\rho_2^j.
             \label{xyf:defectrec}
\end{align}
The defect statements hold at all scales for at most two unit signed
sources, including coincident sources.  For \(k\) fixed sources
separated by at least \(\epsilon_0 n\) in lattice units, they hold
through scales \(L^{j+1}\le c\epsilon_0 n/(k+1)\).
The constants in these decaying recurrences and the required lower
bound on \(b\) are independent of \(k\).  The subsequent bounded
number of steps preserve an \(o(1)\) activity bound as \(n\to\infty\),
with constants allowed to depend on the fixed configuration and \(k\).
\end{proposition}

\begin{proof}
Let \(c_{\rm loc}\) bound, in current coarse-block units, every assignment,
normalization-provenance, and collar neighborhood used by one local map.
This includes inherited coefficient provenance: its earlier radii sum
geometrically in current units, as in the proof of
Proposition~\ref{rg:partition-map}.
Let \(C_X\) bound the diameter of a small input in fine-block units.
The fixed geometric choices include \(H>c_{\rm loc}\) and
\((H-c_{\rm loc})L>H+C_X\), as well as the kernel clearances. Thus a good
output's enlarged hull contains all the support and scalar-normalizer
provenance needed to compute it, with a clearance margin.

An untwisted support is bulk-good if its required neighborhood has
distance greater than \(Hm\) from the physical walls.  On a torus
require coordinate gaps of that size so the neighborhood unwraps.
The remaining supports are wall supports.  Copy the bulk localization
on small good inputs.  On small wall inputs subtract only their
neutral constant; the full untwisted activity is even.  To determine
a good output, every contributing fine support and every assigned
piece lies in its coarse sites or within the fixed one-step
assignment neighborhood.  The clearance ensures that all these
supports are good and that the shell covariance agrees with the
bulk covariance there.  Connectivity tests, temporary exclusions,
and singleton normalizations then agree as well.  This proves
equality with the bulk on good outputs by induction, including the
normalization step.

At a fixed coarse wall output, a small wall input must lie in a strip
of bounded fine-block width along a side.  There are \(O(L)\) such
positions and only a fixed number of small animals at each position.
Corner positions are already included in this count.  Thus even
neutral wall remainders cost \(CL\cdot L^{-2}\), and charged ones
cost \(CL\cdot L^{-4}\).  Small torus-winding supports are absent
in the regime where the torus clearance is used.  Large inputs
contract by the weight as before.  A good input contributing to a
bad output is bounded crudely by \(Cu_j\), including any assigned
tensor polynomial; it is not required to cancel there.  The product
and decoration counts in the proof of Proposition~\ref{xyf:recurrence}
give the quadratic term.  The initial size is \(w_0\le C\eps\),
which proves the first recursion.

For a spin source array \(l=\sum_{r=1}^k\varsigma_r\delta_{x_r}\),
\(\varsigma_r\in\{1,-1\}\), use the exact connection of
Proposition~\ref{xyh:compensation}.  Its offsets are
\[
 \eta=-\mathcal R^{-1}dG_nl/\sqrt b,
 \qquad \vartheta=a_l/\sqrt b,
 \qquad (w,p,\psi)\mapsto
       (w-id_c\eta/\mu,p+\eta,\psi+\vartheta).
\]
The symbol \(a_l\) here denotes the connection lift, not the scalar
\(a=s/d_c\) in the diagnostic.  A support is shift-good if its
rectangular hull enlarged by \(Hm\) misses every marked source.
On that hull, clipped to the physical square when necessary, the
connection has a common lift modulo its period; the reflected
extensions obey the same identity on the stencil neighborhoods.
Let \(\mathcal R_j^{\rm un}(X)\) be the chosen untwisted representation;
each selected live root is untagged. On a good support represent the full
twisted activity by \(T_\Lambda\mathcal R_j^{\rm un}(X)\), where
\(\Lambda=(-i d_c\eta/\mu,\Lambda_\psi)\) and
\(\partial\Lambda_\psi=\eta\) on its enlarged hull.

The first independent-piece integration, temporary grouping, and relocation
of the full linear \(U_e\) expectation construct \(K_0\) before good and
defect outputs are classified. The initial field collar and retained
exclusion provenance lie within a fixed number of unit blocks of each
output support, covered by the chosen \(H\). On an initial good output the
exact microscopic substitution with the common lift commutes with this
finite algebra. The designated contraction scalars are source independent,
while the whole shifted polynomial remains in the activity. Thus
\(K_0^{\rm tw}(X,z)=K_0^{\rm un}(X,z+\Lambda)\) there, and we choose its
representation afresh as \(T_\Lambda\mathcal R_0^{\rm un}(X)\). Initial
defect outputs retain their intrinsic tagged representations, including
shifted \(U_e\) polynomials. The no-return induction below begins at these
classified \(K_0\) outputs and applies only to subsequent \(m\to Lm\) maps.

For localization on small good inputs copy the untwisted pieces and retain their \emph{untwisted} designated
constants.  In particular no background value of a translated
tensor is reassigned as a scalar constant.  On small defect inputs
subtract only the neutral constant.

Every support and normalization neighborhood contributing to a good
output fits within its enlarged hull with a clearance margin. More
explicitly, if \(Y\) is good at \(S=Lm\), every contributing fine
predecessor \(X\), including a predecessor of a scalar normalizer, obeys
\[
 \operatorname{dist}(x_i,\operatorname{Hull}X)
 >(H-c_{\rm loc})S-C_Xm>Hm
\]
for each mark \(x_i\); residual predecessors have the stronger bound
because their closures lie in \(Y\). Thus all fine inputs are good and
admit the same lift. Using the exact fine activity identities with this
common lift, substitution commutes with the convolution, product, and assignment identities.
This proves the good-output identity inductively, including every
designated scalar normalizer and its provenance. After proving that
identity at each still-good output \(Y\), choose its representation
afresh as \(T_{\Lambda_Y}\mathcal R_{j+1}^{\rm un}(Y)\). Thus the
transport is applied to an untagged root; the previously assembled tagged
representation of that good output is discarded. This is a choice of a
representation of the same exact activity.

At a defect output retain, in each selected term, its concatenated live
root with its existing tags and its separately completed coefficient
occurrences. No lift is chosen on that output and no further
\(T_\Lambda\) is applied to that root. An assigned piece originating at a
good fine block is transported only at that fine block, before relocation,
even if its target coarse
singleton is a defect. Its opaque untwisted coefficient is unchanged.
The following no-return property justifies this rule. If a fine support
\(X\) is defective, some mark has distance at most \(Hm\) from its hull.
A residual output contains \(\bar X\), so it is defective at \(S=Lm\).
A singleton receiving a small defect coefficient belongs to \(\bar X\)
and meets \(X\); its distance from that mark is at most \((H+C_X)m\).
Any output whose assignment or normalization neighborhood uses that
singleton is within a further \(c_{\rm loc}S\). The chosen inequality
\((H-c_{\rm loc})L>H+C_X\) makes such an output defective as well.
This includes a defect-specific scalar inverse decoration: it is an
opaque number supported at its singleton, and every later use retains
that support or the stated normalization provenance. Defect localization
emits only a neutral constant, so it cannot create a relocated live
polynomial elsewhere. A good-origin polynomial first landing in a defect
then follows the residual rule. Consequently no tagged defect ancestry,
including its scalar-normalizer provenance, re-enters a good output.

The derivative bounds for the connection give bounded scaled positive
derivatives at every point of a good fine collar for one or two sources.
In the separated \(k\)-source regime, at each folded collar point at most one
source is within \(\epsilon_0n/3\); the identity of this source may vary
with the point. The others contribute at most
\(Ckm/(\epsilon_0n)\) in the scaled bound. The own-birth estimate after
Equation~\eqref{xyh:early-source-scales} applies to every old atom, since
its birth size \(h\le m\) obeys the same early inequality. Choosing
\(c\) small therefore puts the transport of each good input in the
admissible early class.

Apply the explicit comparator in Proposition~\ref{xyh:calculus}(iv) to
that input before relocation:
\[
 A_D^{|X|}\mathsf s_X(T_\Lambda\mathcal R_j^{\rm un})
 =A_D^{|X|}\mathsf s_X(\mathcal R_j^{\rm un})
 \le (A_D/A_U)^{|X|}\|\mathcal R_j^{\rm un}\|_{A_U}.
\]
Its structural coefficient, every opaque completed coefficient size,
and the live event weight are identical on the two sides. All bounded
early contour translations remain inside the endpoint integrand. The
fixed ratio \(A_U/A_D\) is used only in the subsequent spatial
enumeration of support and assignment choices. The product bounds apply
to concatenated live roots and multiply the sizes of completed scalar
occurrences separately. This gives good-input forcing
\(C(u_j+w_j)\), including assigned polynomials on their original good
fine blocks, without a cost for every old event at each scale.

A small defect input has bounded fine diameter and an enlarged hull
meeting a mark.  There are only \(O(1)\) such fine supports per
mark in a coarse output.  In the indicated separated regime the
number of relevant marks in this local calculation is bounded;
for at most two sources it is bounded at every scale.  Consequently
the neutral and charged linear costs are
\(C(L^{-1}+L^{-4})\).  Large defect inputs contract by the weight.
The product counts give the second recursion; the microscopic
offset bounds give \(v_0\le C\eps\).

To solve the recurrences, take the initial sizes sufficiently small
that each quadratic term can absorb a fixed small multiple of the
unknown wall or defect size into a contraction coefficient strictly
less than \(\rho\).  Iterating the resulting linear inequalities
and summing their geometric convolutions gives the displayed
\(C\eps\rho_1^j\) and \(C\eps\rho_2^j\) bounds.  Equivalently,
substitution of these bounds into the two recurrences proves them
inductively after choosing their constants and then increasing \(b\).

For a fixed separated configuration, the last scale in the decaying
regime is proportional to \(n\).  There are only
\(O_{k,\epsilon_0}(1)\) remaining scales, on grids with
\(O_{k,\epsilon_0}(1)\) blocks.  The entering activities tend to zero.
The finite remaining maps and terminal integrals have finite bounds
by the history calculus with the actual source offsets;
their constants may depend on \(k,\epsilon_0,b\).  Each new term has
at least one entering activity, so finitely many applications of
these bounds preserve \(o(1)\). In each live root, the number of atoms
whose own birth sizes are late is bounded by
\(C_{k,\epsilon_0,L}\), as proved after
Equation~\eqref{xyh:early-source-scales}. Only their explicitly recorded
late numerical product may change by a fixed-configuration factor. Early
decorations stay in the uniform endpoint class, and opaque completed
coefficient occurrences are not transported. Thus neither the original
charge and \(F_b\) weights nor the microscopic lower bound on \(b\)
is increased with \(k\).
\end{proof}

Let \(R_n(l)\) be the twisted-to-untwisted ratio of the remaining
compensated partition functions, without the deterministic Gaussian
source factor.  The exact identity of Proposition~\ref{xyh:compensation}
reads
\begin{equation}
 \E_{\mathrm{XY},b,n}e^{i\sum_x l_x\theta_x}
  =\exp\left\{-\frac{s+d_c^2/\mu}{2b}(l,G_nl)\right\}R_n(l).
 \label{xyf:spinfactor}
\end{equation}

\begin{proposition}[A universal spatial amplitude]\label{xyf:amplitude}
There is \(c_b\in(0,\infty)\), depending only on \(b\) and the
fixed surviving parameters, such that, for each fixed \(k\) and
every convergent configuration of distinct interior continuum
locations,
\begin{equation}
 R_n\!\left(\sum_{r=1}^k\varsigma_r\delta_{x_{r,n}}\right)
       \longrightarrow c_b^k.
 \label{xyf:amplitudelimit}
\end{equation}
The same \(c_b\) works for both signs, every interior position,
every integer sequence \(n\to\infty\), and every choice of nested
cuts.  The ratios for at most two sources are bounded uniformly
without a separation condition.
\end{proposition}

\begin{proof}
After extracting the designated constants at each step, the
terminal normalized integrals in both problems tend to one.
There are boundedly many terminal blocks, and
Proposition~\ref{xyf:forced} and the terminal history estimate bound
every nonempty activity contribution by \(o(1)\).  The same reasoning
includes the bounded number of late steps for a fixed separated
configuration.

Write \(1+a_{j,B}^{\mathrm{tw}}\) and
\(1+a_{j,B}^{\mathrm{un}}\) for the extracted factors.
They are real: on real background fields complex conjugation
combined with \(w\mapsto-w\) preserves the compensated integral
and its local rules, including the real connection offsets.
Their distance from one is small in the decaying regime, so they
are positive.  The microscopic designated constants are identical
in the two problems.  At subsequent scales a difference arises
only from a small defect input.  The bounded number of such inputs
per source, the extracted-coefficient estimate, and
Proposition~\ref{xyf:forced} imply
\begin{equation}
 \sum_B\left|\log(1+a_{j,B}^{\mathrm{tw}})
                 -\log(1+a_{j,B}^{\mathrm{un}})\right|
       \le Ck\eps\rho_2^j.
 \label{xyf:logtail}
\end{equation}
This follows from \(|\log(1+x)-\log(1+y)|\le2|x-y|\)
for the small real coefficients in question.  In a separated
configuration the bounded late contributions are \(o(1)\).
For two sources the estimate holds throughout.  The terminal
history bound then gives a uniform ratio bound for all sufficiently
large grids.  The finitely many remaining grid sizes are covered
by enlarging that bound; their positive partition functions are
finite and their source locations form finite sets.

For a single signed mark remaining macroscopically interior,
translate the nested cuts by its lattice position.  This translation
does not change the bulk parameters or the exact physical ratio.
At any fixed depth all constants that differ between the two
problems are computed within a fixed lattice neighborhood of the
mark.  In this neighborhood, as the distance to the boundary tends
to infinity, the grids and kernels are eventually the plane ones.
The connection increments converge to the isolated plane connection.
Its phase values modulo the period converge up to one common phase.
A common phase translates the activity argument in the scalar
field; neutral constant extraction is invariant under that
translation, and the good-input lift rule preserves the invariance
at every step.  The finite-depth continuity of the history integrals
therefore proves convergence of each local constant difference.
Simultaneous negation of the fields and the connection proves that
these limits are unchanged by reversing the source sign.
They are also independent of the macroscopic source position,
because the fixed-depth limiting grid and increments are the same.

The sum of these isolated differences is absolutely convergent by
\eqref{xyf:logtail}.  Define its exponential to be \(c_b\).
It is finite and strictly positive.  The terminal ratio tends to
one, so it is also the limit of the exact single-source physical
remaining ratio.  In particular its value is independent of cuts,
although cuts were translated to compute that limit.

For several separated sources, compare with the individual
single-source problems on the \emph{same} hierarchy.  At each fixed
depth the neighborhoods supporting differing constants around
distinct marks are disjoint for large \(n\).  The other sources'
connection increments tend to zero in each neighborhood; their
phase contribution is a common phase to leading order and thus
does not affect neutral constants.  Finite-depth continuity makes
the local differences converge to the corresponding single-source
differences.  A mark has only finitely many positions modulo the
cuts at a fixed depth, so this comparison is uniform in those
positions.  First sum through a fixed depth, then use
\eqref{xyf:logtail} to send that depth to infinity.  The result is
\(k\log c_b\), proving \eqref{xyf:amplitudelimit} with no
subsequence or location-dependent adjustment.
\end{proof}

\subsection{The full complex spin field}

\begin{theorem}[Low-temperature XY field]\label{xyf:spin}
There is \(b_0(\mathrm{XY})<\infty\) such that for every fixed
\(b\ge b_0(\mathrm{XY})\), the choices
\begin{equation}
 K_{\mathrm{XY}}(b)=\frac{b}{s+d_c^2/\mu},
 \qquad A_{\mathrm{XY}}(b)=c_b^{-1}
 \label{xyf:stiffness}
\end{equation}
satisfy \(K_{\mathrm{XY}}(b)>1/(4\pi)\),
\(A_{\mathrm{XY}}(b)>0\), and \(K_{\mathrm{XY}}(b)/b\to1\).
Along every integer sequence \(n\to\infty\),
\[
 n^{-2}\sum_{x\in D_n^\circ}
  A_{\mathrm{XY}}(b)e^{G_n(x,x)/(2K_{\mathrm{XY}}(b))}
       e^{i\theta_x}\delta_x
       \ \Longrightarrow\ V_{K_{\mathrm{XY}}(b),D}
\]
in \(H^{-3}_{\mathrm{loc}}(D)\), and jointly on any finite
collection of complex smooth compactly supported test functions.
The limit uses the full Dirichlet-variance Wick convention.
\end{theorem}

\begin{proof}
By Proposition~\ref{xyf:shooting},
\(s+d_c^2/\mu=1+O(b^{-1})>0\).  This proves the stated asymptotic
for the coefficient \(K_{\mathrm{XY}}(b)\) and, on increasing \(b_0\),
the required lower bound.
Let \(K=K_{\mathrm{XY}}(b)\), \(A=c_b^{-1}\), and write
\(Z_n(x)=A e^{G_n(x,x)/(2K)}e^{i\theta_x}\).
For signed factors let \(Z_n(x)^{[+]}=Z_n(x)\) and
\(Z_n(x)^{[-]}=\overline{Z_n(x)}\).  Equation~\eqref{xyf:spinfactor}
gives the exact cancellation
\begin{equation}
 \E\prod_{r=1}^k Z_n(x_r)^{[\varsigma_r]}
  =A^k R_n\!\left(\sum_r\varsigma_r\delta_{x_r}\right)
    \exp\left\{-\frac1K\sum_{r<t}
                      \varsigma_r\varsigma_tG_n(x_r,x_t)\right\}.
 \label{xyf:wickcancel}
\end{equation}
In particular every self-energy, including its finite
position-dependent Dirichlet part, cancels.  By
Proposition~\ref{xyf:amplitude} and the lattice Green convergence
in Lemma~\ref{rg:green}, at separated convergent interior points
this tends to
\[
 \exp\left\{-\frac1K\sum_{r<t}
                   \varsigma_r\varsigma_tG_D(z_r,z_t)\right\}.
\]
These are precisely the separated signed moments of the full-variance
Wick exponential.  Convergence along every convergent separated
configuration sequence is uniform on separated compact configuration
sets: otherwise compactness would supply a convergent sequence on
which the error stayed bounded away from zero.

We check the remaining hypotheses of Theorem~\ref{vi:field-transfer}.
The two-source ratio bound and the Green estimates give, uniformly
in all lattice points of a fixed interior compact set,
\begin{equation}
 0\le\E[Z_n(x)\overline{Z_n(y)}]
        \le C(|x-y|+n^{-1})^{-Q},\qquad Q<2.
 \label{xyf:twopoint}
\end{equation}
One may take \(Q=C_G/K<2\), where \(C_G\) is the constant in the
discrete logarithmic Green bound, by increasing the lower bound on
\(b\) once.  Repeated
locations are included: their diagonal Green factor is bounded by
the same estimate.  The same-sign two-point expression is bounded
above by a constant since \(G_n\ge0\).

All signed-spin expectations are nonnegative.  Indeed the Fourier
coefficients of \(e^{b\cos\theta}\) are nonnegative: expand
\(\exp((b/2)e^{i\theta})\exp((b/2)e^{-i\theta})\) and group equal
powers.  Haar integration at interior vertices selects currents
with the required integer divergence.  Every surviving current
has nonnegative weight, and the pinned boundary supplies the
remaining divergence.  The positive deterministic multipliers in
\(Z_n\) preserve this property.

Finally fuse all fixed boundary vertices into one pin vertex and
give that vertex an independent uniform angle, rotating every
spin by the same angle.  The resulting measure is the finite
free ferromagnetic plane-rotator measure, with nonnegative edge
couplings, and this rotation preserves the modulus of each
nonnegative weighted complex smear.  The Lee--Yang estimate in
Lemma~\ref{vi:lee-yang} therefore applies to its real and imaginary
projections and gives the absolute-smear moment bounds needed in
Theorem~\ref{vi:field-transfer}.  Its public theorem hypotheses are
finite graph, nonnegative couplings, and nonnegative smear weights;
all are satisfied here.  General complex test functions are handled
by splitting real and imaginary parts into nonnegative parts, as
in that theorem.

Thus the hypotheses of the transfer theorem are all available:
separated signed multipoint limits for every fixed order, the
unrestricted integrable two-point bound \eqref{xyf:twopoint},
positivity of every signed moment, and the rotated Lee--Yang
moment bound.  That theorem controls collisions, proves moment
determinacy, and gives local negative-Sobolev tightness.  It yields
the claimed full complex law and joint smeared convergence.
The large fixed \(b_0\) was chosen in the bulk and the at-most-two-source
estimates; Proposition~\ref{xyf:forced} supplies every higher fixed
order at this same \(b\).  The arbitrary-size hierarchy and the
cut-independent amplitude prove the assertion for all integers
\(n\), rather than only a geometric subsequence.
\end{proof}
